% Standalone manuscript fragment; no preamble.
% Requires proposition and proof environments (e.g. amsthm).
\paragraph{Why exposure must retain conjunctive guards.}
The following elementary separation explains the conjunction in (GP).

\begin{proposition}[Acquisition exposure is not additive]
\label{prop:acquisition-nonadditivity}
Let $r$ be a retained-record atom, $k$ a wrapping-key atom, and $s$ a
protected-seed atom associated with an origin identifier $o$.
For $A\subseteq\{r,k\}$, let $\operatorname{Cl}(A)$ be the least set
containing $A$ and closed under the rule
\[
  \{r,k\}\;\Longrightarrow\;s.
\]
There are no other rules in this fragment.  Define the exposed-origin set by
\[
  E(A)=\{\,o : s\in\operatorname{Cl}(A)\,\}.
\]
There is no fixed function
$f:\{r,k\}\longrightarrow\mathcal{P}(\{o\})$ satisfying
\[
  E(A)=\bigcup_{x\in A} f(x)
  \qquad\text{for every }A\subseteq\{r,k\}.
\]
\end{proposition}

\begin{proof}
The closure rule gives
\[
  E(\{r\})=E(\{k\})=\varnothing,
  \qquad
  E(\{r,k\})=\{o\}.
\]
Exactness for the two singleton acquisitions forces
$f(r)=f(k)=\varnothing$.  Their union is then empty for the joint
acquisition, contradicting $E(\{r,k\})=\{o\}$.
\end{proof}

\paragraph{Meaning of exposure.}
Membership of an origin $o$ in $E(A)$ records derivability of its
protected seed under the specified closure rules; the identifier itself
may be public. Computational estimation is measured separately by success
probability. A one-time-pad illustration makes the dependence on joint
observations explicit: let $S,W$ be independent uniform $\lambda$-bit strings and
$C=S\mathbin{\oplus}W$, with no additional observation depending on $S$ or
$W$.  Then
\[
  H(S\mid C)=H(S\mid W)=\lambda,
  \qquad H(S\mid C,W)=0.
\]
From either observation alone, any estimator of $S$ succeeds with
probability at most $2^{-\lambda}$; from both, $S=C\mathbin{\oplus}W$
recovers it exactly. The XOR construction serves here as an
information-theoretic illustration of acquisition; the protocol's record
encryption is handled by its cryptographic record interface.

\paragraph{Consequence for provenance routes.}
A source label and an exposure decision serve different purposes.  Copying
a retained record preserves the source label, while subsequent acquisition
of a matching key may change exposure.  A guarded-route representation
retains this distinction: alternative acquired records supply alternative
routes to the same source, and every route retains the conjunction of its
required keys.  Exposure is therefore evaluated from the joint acquisition
and key-derivation history. The separation applies specifically to an
exact exposure representation formed by the union of fixed per-object
labels.
